\section{A uniform symmetric-square subspace of the stable spectrum}
\label{sec:kt-symmetric-square}

The stable spectrum itself contains a uniform classical family.  The
construction in this section uses classical forms of weight $e+1$,
whereas the endoscopic construction used weight $2e+1$.  We identify
the resulting subspace and its exact polynomial factor.  The word
``constructed'' below specifies its scope: we do not classify every
possible adjoint origin over $E$.

\begin{lemma}[KT-GL2-LOWER-SCOPE]
The integral basis, Hecke determinant formula and CM classification
in KT-GL2-CM hold for the parameter $n$ in every integer degree
$n\geq1$, with the empty-space convention when $n=1,2$.
In particular, for even $e\geq2$,
\[
 L_e=S_{e+1}(\Gamma_0(3),\epsilon),\qquad
 \dim L_e=\left\lfloor\frac{e-2}{3}\right\rfloor,
 \qquad
 \dim L_e^{\rm CM}={\bf1}_{6\mid e}.                 \tag{KT-S.1}
\]
Twisting by $\epsilon=\chi_{-3}$ pairs the non-CM normalized
eigenforms without fixed points.
\end{lemma}
\begin{proof}
Use the same graded-ring statement on $\Gamma_1(3)$ as in (GL.1).
The weight and character of $Da^{2n-5-3j}b^j$ give an integral
unitriangular basis with $0\leq j<\lfloor(2n-2)/3\rfloor$;
no part of that graded-ring argument requires $n$ to be even.
For $n=1,2$ the permitted index set is empty.  The Fourier
coefficient formula for $T_p$ has factor $\epsilon(p)p^{2n}$
and is unchanged.  The CM proof uses only the level-three
conductor formula and the unit condition $6\mid2n$, so also
does not require even $n$.  Substituting $n=e/2$ proves (KT-S.1).

For completeness, the local twisting argument is likewise
independent of that parity.  Primitivity of the nebentypus
excludes oldforms.  At three the conductor-one Weil--Deligne
parameter has inertia characters $1,\epsilon$ and monodromy
zero, as in KT-U3-ENDO-INJECT.  Twisting interchanges these
two inertia characters and preserves the conductor.  There
is no change of conductor at the other finite primes.
The twist therefore stays in the same primitive level-three
space.  A fixed point is precisely a CM form for $E$, by the
quadratic self-twist criterion already used there.  Hence the
non-CM forms occur in two-element orbits.
\end{proof}

\begin{theorem}[KT-U3-SYM2-INJECT]
For every even $e\geq2$, the non-CM quadratic-twist orbits
in $L_e$ give distinct stable automorphic eigenlines in
$M_{e-1}$.  Their span $M_{e-1}^{\rm sym}$ has dimension
\[
                 H_e=\left\lfloor\frac{e-2}{6}\right\rfloor.
                                                               \tag{KT-S.2}
\]
The corresponding simple global parameter is
\[
 \Pi_f=\operatorname{BC}_{E/\mathbf Q}
             \bigl(\operatorname{Ad}(\pi_f)\bigr),
 \qquad
 \operatorname{Ad}(\pi_f)
       =\operatorname{Sym}^2(\pi_f)\otimes\epsilon^{-1}.
                                                               \tag{KT-S.3}
\]
At every split prime $p\ne3$, the eigenvalue of the scaled
compact operator $p^{e-1}T_p$ is
\[
                   a_p(f)^2-p^e.                              \tag{KT-S.4}
\]
The same value occurs for the opposite minuscule direction.
\end{theorem}
\begin{proof}
\paragraph{The global lift and its local parameters.}
We use unitary normalizations for the automorphic representations.
The Gelbart--Jacquet theorem, \emph{A relation between automorphic
representations of $\mathrm{GL}(2)$ and $\mathrm{GL}(3)$},
Theorem~9.3, printed p.~534, gives the adjoint lift with its local
parameters at every place and makes it cuspidal when the original
representation has no nontrivial character self-twist.  Its
normalization is the determinant-one adjoint representation:
Definition~3.1.3 and Proposition~3.2, pp.~485--486, identify it
at infinity, and Proposition~3.3, in particular (3.3.6), p.~487,
gives the principal-series local formula used at every finite
place here.  The self-twist case is treated in Section~3.7.
A non-CM holomorphic form here has no such self-twist:
any nontrivial self-twist is quadratic by the central character,
and the quadratic self-twist criterion makes the form CM.
Thus $A_f=\operatorname{Ad}(\pi_f)$ is cuspidal on
$\mathrm{GL}_3/\mathbf Q$, self-dual, and has central character one.

Arthur--Clozel, \emph{Simple algebras, base change, and the
advanced theory of the trace formula}, Chapter~3,
Theorem~4.2(a), p.~202, gives
$\Pi_f=\operatorname{BC}_{E/\mathbf Q}(A_f)$.
Theorem~5.1, pp.~212--213, makes this strong base change at
every finite and infinite place, in the sense of Definition~1.2,
p.~199.  The cuspidality criterion says that failure of
cuspidality would give the
nontrivial self-twist $A_f\simeq A_f\otimes\epsilon$.
The central characters of these two representations are $1$
and $\epsilon^3=\epsilon$, respectively, so this is impossible.
Only quadratic base change and its fibers are used here;
no statement about a general tower of solvable descent is needed.

At every prime other than three, the classical representation
is unramified, so its adjoint and base change are unramified
with zero monodromy.  At three its two-dimensional inertia
representation is $1\oplus\epsilon$ and $N=0$.  The adjoint
inertia is therefore $1\oplus\epsilon\oplus\epsilon$.
Restriction to $W_{E_3}$ is trivial, and monodromy stays zero.
The base change is thus unramified at every finite place of $E$.

At the complex place its characters have exponents
$(e,0,-e)$ in the notation $(z/\bar z)^s$.  This follows either
by applying the adjoint representation to the two real
discrete-series weights, or from the ratios of their complex
restrictions $(z/\bar z)^{\pm e/2}$.
The global representation $\Pi_f$ is both self-dual and invariant
under conjugation.  Its conjugate self-dual parity is positive.
To check the latter with precisely the source convention used
above, apply the global seed theorem and its localization in
KMSW Propositions~1.3.1 and~1.3.3.  At infinity the three
distinct characters $(z/\bar z)^e,1,(z/\bar z)^{-e}$
are each of positive conjugate-dual parity.  Their multiplicities
are one, so a negative conjugate-dual form is impossible there.
The global parity must consequently be positive.

\paragraph{The compact member and multiplicity.}
The representation $\Pi_f$ is a simple generic parameter for
the rank-three unitary group.  Its real packet contains the
compact member of coefficient highest weight $(e-1,0,1-e)$.
At every finite place it has a member fixed by the prescribed
standard compact subgroup: the spherical local correspondence
at split and unramified nonsplit primes, and the checked
ramified principal-series cases of Adler--Lansky at three,
give these members.  Each fixed line is one dimensional,
by KT-U3-LEVEL-CONVERSE and the usual spherical uniqueness.

For a simple generic parameter the reduced global component
group is trivial.  Hence there is no remaining sign constraint
in the generic pure-inner multiplicity formula of KMSW
Theorem~5.0.5.  The tensor product of the compact real member
and these finite members occurs with multiplicity one.  It
is stable because its base change is a single cuspidal
$\mathrm{GL}_3/E$ representation.

\paragraph{Precisely the asserted fibers.}
The adjoint is unchanged under twisting a two-dimensional
representation, so $f$ and $f\otimes\epsilon$ give the same
parameter.  Conversely, suppose $\Pi_f\simeq\Pi_g$.
The fibers of prime-cyclic base change, Arthur--Clozel,
Chapter~3, Theorem~3.1, p.~201 (equivalently the fiber assertion
in Theorem~4.2(d), p.~202), give
$A_g\simeq A_f\otimes\epsilon^j$ for $j=0$ or $1$.
The determinant-one normalization excludes $j=1$, exactly
as in the cuspidality argument.  Therefore $A_f\simeq A_g$.

Ramakrishnan, \emph{Modularity of the Rankin--Selberg $L$-series,
and multiplicity one for $\mathrm{SL}(2)$},
Theorem~4.1.2, printed p.~90, says that equality of the adjoint
representations of two cuspidal $\mathrm{GL}_2$ representations
forces one to be a character twist of the other.  Since the
central characters here agree, the twisting character $\eta$
is quadratic.  At each prime other than three both classical
representations are unramified.  Twisting an unramified
two-dimensional Weil--Deligne representation by a ramified
character makes its inertia nontrivial, so $\eta$ is unramified
outside three.  The only quadratic characters of $\mathbf Q$
unramified outside three are $1$ and $\epsilon$, as follows
from the classification of fundamental quadratic discriminants.
Thus $g=f$ or $f\otimes\epsilon$.

Distinct orbits consequently have distinct simple global
parameters and independent eigenlines.  Their number is
\[
 \frac12\left(
        \left\lfloor\frac{e-2}{3}\right\rfloor
          -{\bf1}_{6\mid e}\right)
        =\left\lfloor\frac{e-2}{6}\right\rfloor.
\]
Uniqueness of the finite fixed lines and global multiplicity
one show that their span has exactly this dimension.

\paragraph{Eigenvalue normalization.}
At a split prime let $\alpha,\beta$ be the classical roots,
so $\alpha+\beta=a_p(f)$ and $\alpha\beta=p^e$.
The adjoint roots are $\alpha/\beta,1,\beta/\alpha$.
The normalization already checked from Gross's Satake formula
multiplies their trace by $p^e$ for $p^{e-1}T_p$.  Thus
\[
 p^e\left(\frac\alpha\beta+1+\frac\beta\alpha\right)
       =\alpha^2+\alpha\beta+\beta^2
       =a_p(f)^2-p^e.
\]
The multiset of roots is invariant under inversion; the
opposite minuscule direction therefore has the same scaled
eigenvalue.  This proves (KT-S.4) for both directions.
\end{proof}

\begin{corollary}[KT-U3-SYM2-POLYNOMIAL]
Let $Q_{e/2,p}(Y)$ denote the classical characteristic polynomial
on $L_e$ for a split $p\ne3$.  Set
\[
 C^{\rm low}_{e,p}(Y)=
 \begin{cases}
 1,&6\nmid e,\\
 Y-(\varpi^e+\bar\varpi^e),&6\mid e,
 \end{cases}
 \qquad\varpi\bar\varpi=p.
\]
Let $H_{e,p}(Y)$ be the unique monic square root of
$Q_{e/2,p}(Y)/C^{\rm low}_{e,p}(Y)$.  Then
\[
 H_{e,p}(Y)\in\mathbf Z[Y],\qquad
 P^{\rm sym}_{e,p}(X)
   =\operatorname{Res}_{Y}
          \bigl(H_{e,p}(Y),X+p^e-Y^2\bigr)\in\mathbf Z[X].
                                                               \tag{KT-S.5}
\]
This is precisely the characteristic polynomial on the
constructed symmetric-square subspace, of degree $H_e$.
There is an exact monic integral factorization
\[
 P^{\rm full}_{e,p}
     =P^{\rm end}_{e,p}P^{\rm sym}_{e,p}P^{\rm rem}_{e,p},
 \qquad
 \deg P^{\rm rem}_{e,p}=D(e-1)-E_e-H_e.              \tag{KT-S.6}
\]
In particular this last degree is
\[
 \begin{cases}
 e^3/216-e/2+2,&e\equiv0\pmod6,\\
 e^3/216-7e/18+20/27,&e\equiv2\pmod6,\\
 e^3/216-7e/18+34/27,&e\equiv4\pmod6.
 \end{cases}
\]
\end{corollary}
\begin{proof}
At a split prime the two members of a quadratic-twist pair
have the same $a_p$.  Removing the unique CM factor therefore
leaves a monic square.  Its monic square root is rational by
coefficient recursion and integral because its roots are
algebraic integers.  The resultant is the product over these
twist orbits of $X+p^e-a_p(f)^2$, so KT-U3-SYM2-INJECT proves
its assertion and shows that it divides the stable polynomial.
The remaining quotient is rational, monic, and has only
algebraic integer roots.  It lies in $\mathbf Z[X]$.
The endoscopic factorization was proved in KT-U3-SPECTRAL-SPLIT.
\end{proof}

\paragraph{Exact independent polynomial consistency.}
The program \texttt{symmetric\_square\_check.py} freshly computes
the lower classical matrices from the uniform integral basis,
for every integer $1\leq n\leq15$.  It checks (KT-S.5)--(KT-S.6)
for all thirty pairs consisting of an even $2\leq e\leq30$
and $p=7,13$.  No archived factorization lists are used.
After the CM removal it takes an exact monic square root,
constructs the companion matrix of $H_{e,p}$, and computes
the characteristic polynomial of $C_{H_{e,p}}^2-p^eI$.
Exact monic division by this polynomial succeeds in every
stable polynomial saved by the earlier audit.  The full
coefficient arrays, fresh lower classical matrices, and
comparison evaluations are saved in
\texttt{symmetric\_square\_check.json}.

At the three new degrees $e=26,28,30$, the symmetric-square
dimension is four, and the remaining stable dimensions are
$72,92,112$, respectively.  For the smaller residual moduli
\[
 697446897311,\qquad323026426255399,\qquad
 236017358500033404991,
\]
the greatest common divisor with the symmetric-square
comparison value at $t_p=1+p^e+p^{2e}$ is the entire modulus,
at both seven and thirteen.  The greatest common divisor with
the remaining stable value is one, at both primes.  Each of
the three larger residual moduli is coprime to both stable
factors at both primes.  These are exact integer gcd assertions;
they require no primality assumption about the listed moduli.

\paragraph{Meaning for a source-aware residual comparison.}
The constructed symmetric-square subspace consists of actual
stable cuspidal parameters.  It cannot be discarded from the
rank-three source comparison on the ground that it is classical
in origin.  If a lower classical form has
\[
               a_p(f)\equiv1+p^e\pmod\lambda
\]
at all split good primes, its two compact minuscule eigenvalues
both satisfy the full Eisenstein trace congruence
\[
       a_p(f)^2-p^e\equiv1+p^e+p^{2e}\pmod\lambda.
\]
More intrinsically, if its residual two-dimensional Galois
representation restricted to $G_E$ has semisimplification
$1\oplus\overline\chi_{\rm cyc}^{\,e}$, then the associated
adjoint restricted to $G_E$ has semisimplification
$\overline\chi_{\rm cyc}^{-e}\oplus1\oplus
 \overline\chi_{\rm cyc}^{e}$.
Thus the all-prime residual condition itself is compatible
with this stable family.  Any elimination argument must
retain its arithmetic source conditions.  The factorization
separates the endoscopic, constructed symmetric-square and
remaining stable obligations, while the original classical
rank-two branch of the comparison is also retained.

\begin{corollary}[KT-U3-SYM2-EISENSTEIN]
Let $e\geq2$ be even and put
\[
      \beta_e=\frac{B_{e+1,\epsilon}}{e+1}.
\]
Suppose $\ell>3$ is a prime for which
$\operatorname{ord}_\ell(\beta_e)>0$.  Suppose also that
either $6\nmid e$ or $\ell\nmid1-2^e$.  Then there is a stable
compact automorphic eigenline in $M_{e-1}^{\rm sym}$ and a
coefficient prime above $\ell$ such that its weight-zero
rank-three Galois representation has residual
semisimplification
\[
 \overline\chi_{\rm cyc}^{-e}\oplus1\oplus
       \overline\chi_{\rm cyc}^{e}
       \quad\hbox{on the whole group }G_E.                     \tag{KT-S.7}
\]
At every split $p\ne3,\ell$, both scaled minuscule eigenvalues
are congruent to $1+p^e+p^{2e}$ at that same coefficient prime.

In particular, this conclusion holds for every prime divisor
$\ell$ of each of the following moduli:
\[
 \begin{array}{c|r}
 e&M_e\\\hline
 26&697446897311\\
 28&323026426255399\\
 30&236017358500033404991.
 \end{array}                                                  \tag{KT-S.8}
\]
No primality assertion about $M_e$ is needed.
\end{corollary}
\begin{proof}
Fretwell--Roberts, \emph{Newform Eisenstein congruences of
local origin}, Theorem~3.6, printed pp.~517--518, applies
with $N=3$, $M=1$, $k=e+1$, $\psi=1$, and
$\phi=\epsilon$.  Its sufficient direction requires only
$k>2$, a coefficient prime away from $6NM$, and positive
valuation of the displayed special-value product.
Here that product is $L(-e,\epsilon)=-\beta_e$.
There is no condition $\ell>k+1$ in this theorem.
It produces a normalized eigenform $f\in L_e$ and one
coefficient prime $\lambda\mid\ell$ with
\[
     a_q(f)\equiv1+\epsilon(q)q^e\pmod\lambda
                 \qquad(q\ne3).
\]
The paragraph immediately after its proof makes the form new
because the nebentypus has full conductor three.

If $6\nmid e$, KT-GL2-LOWER-SCOPE shows that no CM forms
occur in this space.  Otherwise, at the inert prime two a CM
form has $a_2=0$, whereas the constructed form has
$a_2\equiv1-2^e\not\equiv0\pmod\lambda$.
Thus $f$ is non-CM in either case.

The classical Galois representation of $f$, its determinant,
the good-prime congruences, and Chebotarev give
\[
     \overline\rho_{f,\lambda}^{\rm ss}
       \simeq1\oplus\epsilon\overline\chi_{\rm cyc}^{e}.
\]
Restriction to $G_E$ removes $\epsilon$.
Taking the adjoint gives the three displayed composition
factors in (KT-S.7).  Equivalently, use the filtration of the
endomorphisms of a two-dimensional triangular representation
and remove its scalar line; $\ell>3$ makes this removal
unambiguous.  KT-U3-SYM2-INJECT constructs the corresponding
stable compact line.  Compatibility at the good places and
Chebotarev identify its normalized rank-three system with this
adjoint restricted to $G_E$.  Formula (KT-S.4) proves the
claim about both minuscule directions at every split good
prime.

For the numerical cases, the exact generalized Bernoulli
values are
\[
 \begin{aligned}
 \beta_{26}&=\frac{80907883885186609178}{81},\\
 \beta_{28}&=-\frac{516454004244279410402}{3},\\
 \beta_{30}&=\frac{102431533589014497766094}{3}.
 \end{aligned}                                                \tag{KT-S.9}
\]
For each row of (KT-S.8), exact integer arithmetic gives
\[
 M_e\mid\operatorname{num}(\beta_e),\qquad
 \gcd\bigl(M_e,6\operatorname{den}(\beta_e)(1-2^e)\bigr)=1.
\]
Every prime divisor of $M_e$ therefore satisfies all the
preceding hypotheses.  The values in (KT-S.9) can be computed
directly from
\[
 \beta_e=\frac{2\,3^e}{e+1}
                 B_{e+1}\!\left(\frac13\right),
\]
which follows from the defining generalized Bernoulli sum
and $B_{e+1}(1-x)=-B_{e+1}(x)$.
The exact rational values and gcds are recorded in
\texttt{symmetric\_square\_check.json}, with independent
standard-library rational-arithmetic replay in
\texttt{evidence/main\_gate/lower\_bernoulli\_certificate.py}
and its JSON output.  The source-audit sibling independently
read and replayed that certificate.
\end{proof}

\paragraph{Consequence for the higher-degree obstruction.}
The smaller moduli at $e=26,28,30$ are therefore supported by
actual stable automorphic systems satisfying the whole
rank-three residual condition, not just by isolated
determinant coincidences.  No finite or infinite collection
of split spherical Hecke generators can exclude those
systems from the full stable space.  The exact coprimality
of the remaining stable factor locates these particular
obstructions inside the constructed symmetric-square
subspace.  The three larger moduli remain in the original
classical comparison branch.  A successful arithmetic
vanishing proof must use the more specific deformation or
Selmer conditions; the existence of these automorphic
congruences by itself asserts neither realization of an
admissible deformation source nor nonvanishing of a
$K$-group.

\paragraph{Source scope.}
The new construction combines the cited Gelbart--Jacquet,
prime-cyclic Arthur--Clozel, and Ramakrishnan theorems with
the local and pure-inner packet arguments checked in the
preceding sections.  Those published lifting theorems are
imported inputs.  The polynomial computations are consistency
checks of their application and are not a proof of global
functoriality.  No classification of all adjoint origins over
$E$, or vanishing of the symmetric-square arithmetic source
branch, is asserted.
